%% kapa: exam=1 round=26 year=2015 period=1 form=A subject=civil_law subject_ko=민법 seq=17 number=17
%% source: Q-Net 감정평가사 기출문제 (cst00309) artlSeq=5138001 file=extracted/1008/5138001/제26회_감정평가사 1차_기출/제26회_감정평가사 1차_기출/1교시_감평_A형.pdf
%% conversion: kordoc 4.24.51 -> tools/gen1.py (auto); review=proofread-sonnet; flags=none
\begin{question}{17}
기간의 계산에 관한 설명으로 옳지 않은 것은?
\begin{choices}
  \item 2015년 6월 16일 오후 3시부터 10일간이라고 하면, 2015년 6월 26일(금)24시에 기간이 만료한다.
  \item 2015년 4월 1일 오전 10시부터 6개월간이라고 하면, 2015년 10월 1일(목)24시에 기간이 만료한다.
  \item 2015년 10월 1일 오전 0시부터 3개월간이라고 하면, 2015년 12월 31일(목)24시에 기간이 만료한다.
  \item 2015년 6월 28일 오전 10시에 출생한 아이는 2034년 6월 27일(화)24시에 성년이 된다.
  \item 정관에 달리 정함이 없다면, 2015년 4월 29일에 사단법인의 사원총회를 개최하기 위해서는 2015년 4월 22일(수)24시까지 소집통지를 발송하여야 한다.
\end{choices}
\end{question}
\answer{5}
